\subsection{一直在等待的人}

周被小夜送回家，在打开家门的一瞬间，真昼就扑了上来，差点让他一屁股摔在身后的走廊上。\\

亚麻色的面纱以流水之势紧紧抱住了自己，突然受到冲击的周踉跄几步，才总算稳住了身子没有摔倒。然后他关上家门，轻轻将手环绕在她微微颤抖的娇小后背上。\\

让她等了太久了。从真昼的角度来看，在周被小夜带走的那一刻起，她一定就不安得无可奈何。更何况周一直没有回来，她脑海中想必浮现出了一两个不好的想象吧。\\

「我回来了。抱歉回来晚了」

「……你没被怎么样吧？」\\

这在真昼身上真的是非常罕见，没有进门时的寒暄，也没有慰劳的话语，只是最先开口关心周的安危，看来她确实非常不安。\\

「你这也太怀疑了吧。……她没有做任何会让你担心的事，甚至还请我吃了一顿饭。我好好地跟她谈过了」

「……是吗」\\

虽然周说完全没有发生真昼所担忧的那种事，但真昼看起来似乎并没有多信服。想必小夜的所作所为在真昼的记忆中深深扎下了根。\\

（这种地方一般人是会记恨的哦）\\

虽然事到如今再说这些也没用了，但周还是在心里轻轻嘀咕了一句抱怨的话，同时像是在安抚一般，温柔地拍了拍她那感觉比平时要小得多的后背。\\

「抱歉让你操心了」

「不是的。这都怪我把事情全都推给了周君……那个，本来，这是我必须要去面对的事情」

「可是这对真昼来说很痛苦对吧？」

「话是这么说没错」

「既然我能代替，那我觉得由我来代替也是可以的」\\

反过来说，在周看来，有些事情并不一定非得本人去承受。\\

他不认为无论什么事情知道真相就一定是好事。\\

与其被极其痛苦的现实击垮内心，有时用幸福的谎言将其掩盖，给残酷的现实加上一层滤镜，不也是必要的吗？有时甚至可以故意不去看，从而让人能够从眼前的黑暗中转向光明。\\

就算一直等到伤口结痂并干净利落地脱落的那一刻，又有什么不可以呢？\\

「确实，考虑到将来的事，去面对并跨越它才是为了真昼好。但是，如果真昼现在觉得痛苦，我认为没必要勉强自己去面对。在没有准备好的状态下去面对，大概会摔倒的」

「……周君对我太温柔了，也太娇纵我了」

「是真昼对自己太严厉了啦。……如果是受了重伤的人，真昼能在伤口还没愈合的情况下，叫他马上回工作岗位吗？」

「说不出口，我会想让他好好休养」

「那就是这么回事了」\\

心地善良的真昼是不会对别人提出这种无理要求的。同样，别人也会对她这么说。\\

「心里的伤，不是别人能衡量出来的。真昼只要在自己想面对的时候去面对就好了。而且很幸运，这并不会给你今后带来什么不好的影响」

「……嗯」\\

虽然似乎还有话想说，但既然这也是真昼不愿听的事，而且周也不打算主动说，于是真昼在被周说服的形式下点了点头。\\

正如告诉她的那样，就算不听对真昼也不会有什么影响，今后她们也不会对真昼做什么。是好也罢是坏也罢，她都已经真正地置身事外了。\\

「在你在意的范围内，在不需要勉强才能听的范围内，我可以告诉你。你想听吗？」\\

话虽如此，什么都不听对真昼的心理健康可能也不太好，于是周小心翼翼地试探着问道。真昼稍微犹豫了一下，然后垂下了长长的睫毛。\\

「……那个人的优先顺序，排在前面的是慧君，这一点没错吧」

「是啊。……关于这一点，我觉得你有权去怪罪。即使那两人有苦衷，但受到牵连的真昼，有权利去怪罪你的父母和慧君的父亲」

「毫无疑问，肯定是有什么苦衷的，对吧」

「……大概吧」

「既然如此，我也无能为力了，而且，已经无所谓了。都这时候了，我也不可能再对那些人抱有什么期待」

「这样啊」\\

不知是因为周没有详细说明，还是因为知道周不会再作进一步解释，真昼并没有强求周说明之外的事情，很干脆地接受了。\\

那份冷淡反而让人感到心痛，所以周温柔地随声附和，就那样让真昼抱着走上了走廊。\\

大概是在周开锁的同时她就冲到了门口，走廊上乱糟糟地落着一条毯子，就好像真昼蜕下的皮似的，真昼慌忙将它捡了起来。一想到让她等了那么久，周心里充满了歉意。他们收拾妥当来到客厅，在沙发上稍事休息。\\

在周被小夜缠住的时候，也许是为了转移注意力，矮桌上散落着参考书和杂志。而且，那只放在周房间里的、大概被当成了周的替身的猫咪玩偶，正占据着沙发。\\

沙发附近真的有一种前所未有的凌乱感，这也是真昼感到不安的证明。\\

「……觉得稍微清爽了一些」\\

抱着那只大概一直注视着她不安模样的玩偶，真昼小声嘟囔道。\\

「清爽？」

「可能说接受了比较贴切吧。……有一种再次认同了的感觉，那些人果然还是那样的人」

「……嗯」

「看周君的表情，我也知道，在周君看来，这其中隐藏着相当重大的事情，而且那作为他们那么做的理由，在某种程度上也是能让人接受的」\\

周虽然没有详细说明，但也许是因为周保持着平淡的态度，所以她察觉到小夜她们那边并没有说出极其不讲理的话吧。

现状是对方算是一个能讲得通理的沟通对象，虽然从周优先考虑真昼的情感上来说无法接受，但了解了她们的背景后，虽然不能原谅，但也明白了她们那样做的理由。目前大概就是这样的状态。\\

「对我来说，那些人的苦衷我才不管，但看来确实是有正当的理由，而且他们是能够优先考虑那个理由的人，让我再次释然了」

「真昼……」

「那个人看起来也不像是只会凭感情行动的人，是一个能用理性控制感情、比起人情更重逻辑的人。……虽然她是个过分的人这一点并没有改变」\\

似乎在很久以前就已经对期待感到疲惫的真昼，对于自己被轻视的事实，也只是无力地笑了笑就过去了。\\

「而且，如果对那些人来说，我是不被期望出生的孩子的话，那也是没办法的事。我想我在他们眼里大概连放在天平上衡量的价值都没有吧」\\

这种话，能让本人说出口吗。

不对，让她说出这种话的，是小夜他们，也是周。\\

明明没有哭泣的样子，但不知为何看起来却像是要哭出来一般。看着浮现出软弱笑容的真昼，周不由自主地伸出手，将她娇小的身体包覆住，紧紧地抱在了怀里。\\

「……就算对那些人来说是这样，就算真昼也是这么认为的」

「周君……？」

「那个，虽然我觉得，不能因为我期望如此就让你痛苦的过去一笔勾销。……但是，我期望着你的存在，也不打算放手。对我来说，最重要的是你，最优先的也是你。……我，绝对，不会放开你，也不会离开你」\\

真昼，看错周对她的感情了。

周已经在心里决定非真昼不可，即使知道了她的境遇和父母的苦衷，他依然打算连同她的人生一起背负。只要真昼继续渴求周，周就不会放开真昼。\\

即使人心易变，周也认为那是从恋情转变为爱情，然后再转变为家人的亲情。而且，随时都能重新孕育出爱情。\\

周看着自己父母的身影。他们无数次发现对方的优点、尊敬对方并将其转化为爱情。直到现在，他们依然在增加爱情的总量的同时，和睦地作为夫妻生活着。\\

周希望，他们也能成为那样的形式。\\

周凝视着真昼的脸庞，缓慢地诉说着，眼前那双焦糖色眼眸的轮廓微微晃动，仿佛晕染开来一般。\\

「而且，就算真昼不担心，我这人也是超级黏人的。我有绝对不会放手的自信哦？」

「明明是我比较黏人？」\\

周带着十足的调皮笑了起来，真昼也开怀地笑了出来，把脸埋进了周的胸口。\\

「真昼这是在小看我的爱」

「……我才没小看呢。倒是周君是不是小看我的沉重了？」

「真昼的重量我可是很清楚的」

「不是在说物理上的事吧」

「两边我都很清楚」

「……笨蛋」

「我就是笨蛋啊」\\

虽然说法有些不好听，但这半年来周深切地感受到，真昼虽然独立，但执着心其实相当强。在身体方面，真昼有时会靠在周身上，偶尔也会开玩笑地压过来，所以她大概有多重周是知道的，而且他觉得那是一份幸福的重量。\\

周觉得她其实可以增加一下自己在这个世界上的占比，但她似乎有自己的倔强和审美意识，所以只要她自己觉得满意就好。\\

「……又会变得沉重的哦，可以吗？」

「为了能承受住真昼的重量，我平时可是有在刻苦努力的」

「你这不就是在说很重吗」

「倒不如说因为很轻，为了不让我跑到别的地方去，你就算踩着我也没关系哦？为了不让我被吹飞，请好好地作为楔子把我固定住吧」

「才不会被吹飞呢，明明就算被推开，也会自己急急忙忙跑回来的」

「你很了解嘛」\\

付出了那么多的努力才获得了待在真昼身边的权利，周怎么可能主动放弃那个权利。像现在这样紧紧抱着正依赖着自己的真昼，他可没打算放手，而且在真昼说讨厌之前，他都打算待在这里。\\

如果别人想要插队到周的位置，他一定会堂堂正正地把对方赶走。\\

周就是如此信赖自己堆砌起来的努力之山，那已经切实地化作了周身上的自信。\\

看到周堂堂正正地点头，真昼动了动身子，把脸压在周那比以前稍微结实了一点的胸膛上。\\

「……其实周君那么好的人，我这种哪配得上啊」

「真昼，你再说一次，我就堵住你的嘴一分钟哦」\\

交往前她经常把周说得多么自卑之类的，但周觉得真昼才是有过之而无不及。真昼付出的血汗，周和小雪都很清楚。咬紧牙关努力得来的成果，周不想让她用「这种」来形容。\\

这是真昼不好的地方，周抓住她的肩膀，温柔地将她的脸拉开，直直地注视着她。真昼用力地眨了几次眼睛后，露出了一个软乎乎的、近乎哭笑不得的表情。\\

「周君是让我骄傲的恋人。……我一直想着要配得上你」

「这话该让我说才对吧。……对我来说，真昼是最棒的女朋友。懂了吗？」

「……嗯」

「很好」\\

倒不如说，如果不是真昼根本不可能满足，周就是被真昼迷到了这种程度，而且也已经习惯了真昼。这可不是能用「这种」来形容的，希望她能有更多的自知之明。\\

看到周满意地笑了起来，真昼也像被传染了一样，浮现出明朗的笑容。\\

「……顺便一提」

「顺便一提？」

「不说的话，就不帮我堵住了吗」

「你想让我堵住吗？」\\

如果真昼期望的话，真昼想要多少热量，周就愿意分出去多少。但在周近距离凝视着真昼的脸的那一刻，真昼就慌慌张张地用自己的手堵住了周的嘴唇。\\

「……现、现在不行」

「真遗憾」

「真是的」\\

虽然是开玩笑，但他确实喜欢真昼到了有机会就想亲下去的程度，但他也希望真昼不要知道周的这份纠葛。\\

「……那个，我并没有像周君担心的那样消沉。只是感觉了然了。那个，实际见到那个人的时候虽然心悸了一下，但现在已经平静下来了」

「……真的？」

「真的。请相信我」\\

「要确认一下吗？」说完，真昼又带着挑衅的语气在周耳边低语道，看来正如她本人所说，并没有受到太大的伤害，周这才松了一口气。\\

对于那份挑衅，周戳了戳真昼的脸颊作为惩罚，见真昼没有逃跑，周干脆伸手揽住真昼的膝弯，将她抱坐在自己膝盖上重新抱紧。\\

虽然有些惊讶，但真昼果然没有逃跑。反倒如刚才所说，伴随着一声克制的「嘿」，将体重稳稳地压了过来，周也笑着享受这份没多少分量的压迫感。\\

「比起过去的事，我痛切地感受到，果然周君才是最优先的」

「我？」

「对我来说，比起那些人占据内心的比例，周君一个人占据的比例要大得多。简直就像是我的大部分都是由周君组成的一样」

「那确实很高兴，但也把千岁他们加进去吧」

「呵呵，他们和周君是分开算的，也是重要的人。……是啊，我来到这里之后，重要的人和事物比想象中增加了不少呢」\\

曾经的真昼，没有交任何亲密的朋友，只是不断地扮演着被理想化的「好孩子」。没有重要的事物，也没有重要的人。

而现在，真昼开始向许多不想放手的事物伸出手了。\\

「我觉得这是好事。……只要能承受得住，增加也没关系。如果拿不下了，我们再把不需要的东西分类整理掉。……等你能把过去的事情当作是已经理清楚的东西了，我再把今天谈话的内容告诉你」\\

如果拿不下了，把不需要的东西从手中放下就好。只留下回忆，放手就好。对需要的东西进行取舍，并不是什么都要背负。\\

周打算悠闲地等到那粘附在胸口深处、由心中的血凝结而成的痂自然脱落的一刻。\\

「……在那之前能帮我保管我很高兴，但那样周君的容量就会减少了。我有点，讨厌那样。如果是满满的都是我，我会更高兴的」

「独占欲出现了！」

「没错，就是独占欲，也是嫉妒」

「好高兴」\\

基本上很内敛、不太会任性、尽量不表现出来的真昼，用自己的话语表达了对周的执着。这说明对她来说，周的存在占据了相当大的比重。\\

真昼说到底是个内敛且尊重周的人，所以从来不会优先考虑自己。周非常希望她能多依靠他，多独占他。\\

「……周君的这种地方，既是优点也是缺点。会让我得意忘形的」

「你完全可以得意忘形哦？这样我娇纵你的余地就更多了」

「就是这种地方」\\

最近周才知道，真昼的「就是这种地方」这句话，很多时候是在掩饰害羞。表面上听起来像是在挑毛病，但其实心里并不讨厌，实则是欢迎的。周觉得这么理解他自己更高兴，所以就当是这样了。\\

看着努力装出一副冷傲模样的真昼，周笑着再次用力，抱紧了发出轻哼声的真昼的后背，尽情享受她可爱的撒娇。
